% id: 89-13
% book: 베이직쎈 확률과 통계 · 05 확률변수와 확률분포
% section: 유형 051 확률질량함수의 성질; 확률질량함수가 주어진 경우
% figure: none
% answer: 
% answer_source: 
% variant_level: 0 (원본 전사)
% 이 파일은 build.mjs 가 items.json 에서 만든다. 고칠 때는 items.json 을 고친다.
\smash{\raisebox{4.5mm}{\dmkichul{베이직쎈 확률과 통계 89쪽 13번}}}
확률변수 $X$의 확률질량함수가\par\smallskip\dispeq{$\mathrm{P}(X=x)=\dfrac{x+2}{k}\ (x=0{,}\mkern3mu\mathopen{}1{,}\mkern3mu\mathopen{}2)$}\par\smallskip\noindent 일 때, 상수 $k$의 값은?
\dmchoicesauto{}{$5$}{$6$}{$7$}{$8$}{$9$}
